\section{What's Next: Benchmark thế hệ mới và benchmark gần thực tế}
\label{sec:whats_next}

Hai phần trước đã chỉ ra hai giới hạn cốt lõi của benchmark hiện hành: chúng đo
một \emph{lát cắt hẹp} của năng lực và viện dẫn nó như đại diện cho ``tiến bộ AI''
nói chung~\cite{everything_benchmark}, đồng thời chúng \emph{bão hoà} nhanh khi mô
hình mạnh lên, khiến một con số tổng không còn phân biệt được các hệ thống ở
đỉnh~\cite{superglue}. Phần này trả lời câu hỏi \emph{What's Next}: cộng
đồng đang dịch chuyển benchmark theo hướng nào để khắc phục hai giới hạn đó.

\subsection{Từ trắc nghiệm tĩnh sang tác vụ gần thực tế}

Thế hệ benchmark cổ điển (GLUE, SuperGLUE, MMLU) phần lớn là \emph{tác vụ tĩnh,
đầu ra ngắn}: phân loại câu, chọn đáp án trắc nghiệm, điền nhãn. Định dạng này dễ
chấm tự động nhưng cách xa cách con người thực sự dùng mô hình. Một mô hình trả
lời đúng câu hỏi MMLU về ``software engineering'' không đảm bảo nó \emph{sửa được
một bug thật} trong một kho mã thật. Khoảng cách giữa \emph{construct mục tiêu}
(năng lực kỹ sư phần mềm) và \emph{proxy task} (trắc nghiệm về kỹ sư phần mềm)
chính là vấn đề construct validity mà~\cite{measuring_what_matters} cảnh báo.

Benchmark thế hệ mới thu hẹp khoảng cách này bằng cách lấy \emph{chính tác vụ thực
tế} làm đề thi: giải issue GitHub có thật (SWE-bench), hoàn thành một quy trình
nhiều bước trong môi trường thật, hay thao tác trên trình duyệt/hệ điều hành. Đặc
trưng chung:

\begin{itemize}
  \item \textbf{Đầu vào lấy từ thực tế}, không do người ra đề bịa ra, nên khó bị
    tối ưu hẹp và phản ánh đúng phân phối công việc thật.
  \item \textbf{Chấm bằng thực thi} (execution-based) thay vì so chuỗi: chạy thử
    kết quả của mô hình trong một môi trường và kiểm tra tiêu chí pass/fail khách
    quan, thay vì khớp văn bản với một đáp án mẫu.
  \item \textbf{Đầu ra dài, nhiều bước}: lời giải là một patch, một chuỗi hành
    động, hay một phiên làm việc --- không phải một nhãn duy nhất.
\end{itemize}

\subsection{Agentic benchmark: đánh giá cả quá trình, không chỉ câu trả lời}

Khi mô hình được trang bị công cụ (chạy code, đọc file, tìm kiếm, gọi API) và
hành động qua nhiều vòng, ta gọi nó là một \emph{agent}. Đánh giá agent đòi hỏi
benchmark phải cung cấp một \emph{môi trường} chứ không chỉ một câu hỏi: agent
phải quan sát trạng thái, ra hành động, nhận phản hồi, rồi lặp lại cho tới khi
hoàn thành mục tiêu. Điều này kéo theo vài hệ quả về đo lường:

\begin{itemize}
  \item Tiêu chí thành công là \emph{trạng thái cuối} của môi trường (test
    suite xanh, file đúng định dạng), không phải một câu trả lời tức thời.
  \item Cần kiểm soát \emph{tài nguyên} (số bước, số lần gọi công cụ, thời gian)
    vì agent có thể ``đi vòng'' rất lâu.
  \item Xuất hiện rủi ro đánh giá mới: rò rỉ lời giải qua môi trường, hành vi
    không xác định (non-determinism) do công cụ bên ngoài, và chi phí chạy cao.
\end{itemize}

Hai dịch chuyển này --- \emph{gần thực tế} và \emph{agentic} --- là mạch ``What's
Next'' của tutorial. Phần~\ref{sec:swe_bench} lấy \textbf{SWE-bench} làm case study
tiêu biểu, và demo của nhóm (Phần~\ref{sec:demo_design}--\ref{sec:repro})
mượn đúng tinh thần ``chấm bằng cơ chế trích xuất/thực thi'' để cho thấy: khi
benchmark tiến gần thực tế, \emph{chỗ dễ sai} của phép đo dịch từ ``câu hỏi'' sang
\emph{``cách chấm''}.
